\documentclass[12pt,a4paper]{article}
\usepackage[margin=25mm]{geometry}
\usepackage[T1]{fontenc}
\usepackage[utf8]{inputenc}
\usepackage{lmodern}
\usepackage{graphicx}
\usepackage{amsmath,amssymb}
\usepackage{booktabs}
\usepackage{tabularx}
\usepackage{float}
\usepackage{subcaption}
\usepackage{hyperref}
\usepackage{fancyhdr}
\usepackage{titlesec}
\usepackage{listings}
\usepackage{xcolor}

% Code listing style
\definecolor{codegreen}{rgb}{0,0.6,0}
\definecolor{codegray}{rgb}{0.5,0.5,0.5}
\definecolor{codepurple}{rgb}{0.58,0,0.82}
\definecolor{backcolour}{rgb}{0.95,0.95,0.92}

\lstdefinestyle{cmdstyle}{
    backgroundcolor=\color{backcolour},   
    commentstyle=\color{codegreen},
    keywordstyle=\color{blue},
    numberstyle=\tiny\color{codegray},
    stringstyle=\color{codepurple},
    basicstyle=\ttfamily\small,
    breakatwhitespace=false,         
    breaklines=true,                 
    captionpos=b,                    
    keepspaces=true,                 
    numbers=none,                    
    showspaces=false,                
    showstringspaces=false,
    showtabs=false,                  
    tabsize=2
}
\lstset{style=cmdstyle}

\graphicspath{{./}{figures/}{../EMG practicle/}}
\hypersetup{colorlinks=true,linkcolor=black,urlcolor=blue}
\pagestyle{fancy}
\fancyhf{}
\lhead{Biomechanics}
\rhead{Motion Analysis Practical}
\cfoot{\thepage}
\setlength{\headheight}{15pt}
\titleformat{\section}{\Large\bfseries}{\thesection}{0.7em}{}
\titleformat{\subsection}{\large\bfseries}{\thesubsection}{0.7em}{}
\newcommand{\csvfile}[1]{\texttt{\detokenize{#1}}}

\begin{document}

\begin{titlepage}
  \centering
  \vspace*{7mm}
  \includegraphics[width=4.2cm]{campus_logo.png}\par\vspace{7mm}
  \vspace{2mm}{\large University of Moratuwa\par}
  \vspace{18mm}{\large B.Sc. Engineering (Hons)\par}
  \vspace{5mm}{\Large\bfseries Department of Mechanical Engineering\par}
  \vspace{15mm}{\Huge\bfseries 3D Motion Analysis Practical Report\par}
  \vspace{5mm}{\LARGE C3D Binary Motion Capture Parsing, Kinematic Analysis, and Web-Based Simulation\par}
  \vfill
  \renewcommand{\arraystretch}{1.35}
  \begin{tabular}{rl}
    \textbf{Module Code:} & BM3500 \\
    \textbf{Module Title:} & Biomedical Engineering Applications \\
    \textbf{Laboratory:} & Bionics Laboratory \\
    \textbf{Student Name \& Index:} & Nuwanaka W.A.S. (230449N) \\
    \textbf{Facilitators:} & Ms. E.M.A.I. Ekanayake, Mrs. N.B.N. Nuwarapaksha \\
    \textbf{Date:} & 20/10/2026 \\
  \end{tabular}
  \vfill
  {\large \today\par}
\end{titlepage}

\tableofcontents
\newpage

\section{Introduction}

Three-dimensional optical motion capture (mocap) represents the gold standard in biomechanical kinematics, gait analysis, orthopaedic evaluation, and sports rehabilitation. Optical tracking utilizes calibrated multi-camera stereophotogrammetry to triangulate the spatial coordinates ($X, Y, Z$) of retroreflective passive markers affixed to identified anatomical landmarks. By observing marker trajectories across successive time frames, biomechanists extract lower-limb joint centers, segment orientations, linear and angular derivatives, and joint kinematics.

Data collected during optical motion capture sessions are commonly stored in the binary Coordinate 3D (\textbf{C3D}) format. Developed by Andrew Dainis in 1987 and standardized by Motion Lab Systems, the C3D specification encapsulates 3D marker coordinate streams, synchronized multi-channel analog signals (e.g., electromyography and ground reaction forces), and extensive experiment metadata within a unified binary structure. Unlike unstructured comma-separated or tab-delimited ASCII tables, C3D files maintain rigorous parameter blocks, bit-level scaling factors, and camera residual metrics that ensure high numerical precision and compact file sizes.

The primary objective of this laboratory practical is to process, analyse, and visualize lower-extremity gait capture data stored in C3D binary records (\csvfile{Group2_2026.c3d} and \csvfile{Motion2026.c3d}) and ASCII coordinate exports (\csvfile{Group2_2026_(Coordinates).txt} and \csvfile{Group2_2026.txt}). The practical pipeline encompasses:
\begin{enumerate}
    \item Parsing the C3D binary file structure, extracting header parameters, frame counts, sampling rates, and 3D coordinate trajectories.
    \item Reconstructing spatial trajectories for four key anatomical markers: the left hip, left knee, left ankle-bone, and left foot.
    \item Computing linear kinematic derivatives (velocity and acceleration) via numerical finite difference approximations.
    \item Deriving dynamic lower-limb joint angles (knee flexion-extension and ankle dorsiflexion-plantarflexion) using three-dimensional vector dot product formulations.
    \item Simulating the biomechanical motion sequence interactively within a custom WebGL/Three.js rendering engine (\csvfile{simulator.html}) supporting real-time kinematic telemetry.
\end{enumerate}

\section{Methods}

\subsection{Motion Capture System \& Experimental Setup}
The experimental trial was recorded in the Bionics Laboratory using a multi-camera optical motion capture system equipped with stroboscopic infrared (IR) light-emitting diode (LED) rings. Retroreflective passive spherical markers were placed directly on anatomical landmarks of the subject's left lower extremity following a modified Helen Hayes / Plug-in Gait marker protocol:
\begin{enumerate}
    \item \textbf{Left Hip Marker ($H$):} Positioned over the lateral aspect of the greater trochanter / anterior superior iliac spine (ASIS) to approximate the pelvic and hip joint center.
    \item \textbf{Left Knee Marker ($K$):} Affixed to the lateral femoral epicondyle, defining the flexion-extension axis of rotation for the tibiofemoral joint.
    \item \textbf{Left Ankle-Bone Marker ($A$):} Placed on the lateral prominence of the lateral malleolus, identifying the talocrural joint axis.
    \item \textbf{Left Foot Marker ($F$):} Positioned on the dorsal aspect of the fifth metatarsal head / calcaneus to capture foot progression and ground contact events.
\end{enumerate}

The motion capture volume was calibrated using a dynamic calibration wand and an L-frame origin corner, establishing an orthogonal right-handed global Cartesian coordinate system $(X, Y, Z)$ where:
\begin{itemize}
    \item $+X$ represents the anterior-posterior axis (direction of forward gait progression),
    \item $+Y$ denotes the medial-lateral axis (perpendicular to progression),
    \item $+Z$ designates the vertical axis (aligned against gravity).
\end{itemize}
The optical sampling frequency was set to $f_s = 100\,\mathrm{Hz}$, yielding a uniform sampling period of $\Delta t = 0.010\,\mathrm{s}$.

\subsection{C3D Binary File Architecture}
The C3D format is organized into consecutive 512-byte binary blocks, partitioned into three principal sections as outlined in Table~\ref{tab:c3d-structure}:
\begin{enumerate}
    \item \textbf{Header Record (Block 1):} Contains critical file allocation pointers and core acquisition parameters encoded as 16-bit words and 32-bit floating-point numbers. Byte 1 specifies the block pointer where the Parameter Section begins (typically Block 2), while Byte 2 designates the target processor architecture (\texttt{0x50} for Intel little-endian, \texttt{0x54} for DEC, or \texttt{0x55} for MIPS). Word 2 stores the number of active 3D trajectories ($N_{\text{points}}$), Word 3 indicates analog channels per 3D frame, Words 4--5 define the start and end frame indices, and Words 11--12 store the optical camera frame rate in Hertz.
    \item \textbf{Parameter Section (Blocks $2 \dots N_{\text{data}}-1$):} Stores grouped metadata tables. Groups such as \texttt{POINT}, \texttt{ANALOG}, and \texttt{TRIAL} define marker labels (\texttt{POINT:LABELS}), sampling rates (\texttt{POINT:RATE}), trajectory counts (\texttt{POINT:USED}), total frames (\texttt{POINT:FRAMES}), and 3D scale factors (\texttt{POINT:SCALE}). If \texttt{POINT:SCALE} is negative, point coordinates are formatted as IEEE 754 32-bit floating-point values; if positive, coordinates are stored as signed 16-bit integers scaled by the scale factor.
    \item \textbf{Data Section (Blocks $N_{\text{data}} \dots N_{\text{end}}$):} Contains sequential frame records. For each frame, every marker entry occupies 16 bytes (four 32-bit floats $[X, Y, Z, R]$ or four 16-bit words). The fourth component $R$ contains camera contribution bitmasks and reconstruction residuals (reprojection errors). Marker values recorded in millimetres are converted to SI units ($\mathrm{m}$).
\end{enumerate}

\begin{table}[H]\centering
\caption{Structural breakdown of the C3D binary file format.}
\begin{tabularx}{\textwidth}{llX}
\toprule
\textbf{Section} & \textbf{Byte Offset / Blocks} & \textbf{Key Functional Parameters} \\
\midrule
Header & Block 1 (Bytes 1--512) & Parameter pointer, processor ID (\texttt{0x50}), marker count, frame bounds, frame rate ($100\,\mathrm{Hz}$), data start block \\
Parameters & Blocks 2 to $N_{\text{data}}-1$ & Parameter groups (\texttt{POINT}, \texttt{ANALOG}); marker labels, channel mappings, dimension metadata \\
3D Point Data & Block $N_{\text{data}}$ onward & Sequential frame records ($X, Y, Z$ positions, camera residual bitmask) \\
Analog Data & Interleaved with points & Synchronized EMG and force transducer voltage measurements \\
\bottomrule
\end{tabularx}
\label{tab:c3d-structure}
\end{table}

\subsection{Data Processing Pipeline}
The mathematical extraction and derivation pipeline translates raw coordinate streams into kinematic metrics:

\subsubsection{Coordinate Extraction and SI Unit Conversion}
Raw point coordinates $[X_{\text{raw}}, Y_{\text{raw}}, Z_{\text{raw}}]$ are read sequentially across all valid frames $n \in \{1, \dots, N\}$. Coordinates stored in millimetres are scaled to SI units:
\begin{equation}
\vec{r}_m[n] = \begin{bmatrix} X_m[n] \\ Y_m[n] \\ Z_m[n] \end{bmatrix} = \frac{1}{1000} \begin{bmatrix} X_{\text{raw}}[n] \\ Y_{\text{raw}}[n] \\ Z_{\text{raw}}[n] \end{bmatrix}, \quad m \in \{H, K, A, F\}.
\end{equation}

\subsubsection{Linear Velocity and Acceleration Derivation}
Linear velocity vectors $\vec{v}[n] = [v_x[n], v_y[n], v_z[n]]^T$ are determined using a second-order central difference scheme for interior time steps ($1 < n < N$):
\begin{equation}
v_x[n] = \frac{X[n+1] - X[n-1]}{2\Delta t}, \quad v_y[n] = \frac{Y[n+1] - Y[n-1]}{2\Delta t}, \quad v_z[n] = \frac{Z[n+1] - Z[n-1]}{2\Delta t}.
\label{eq:central-vel}
\end{equation}
Linear acceleration vectors $\vec{a}[n] = [a_x[n], a_y[n], a_z[n]]^T$ are similarly computed via second-order central differences:
\begin{equation}
a_x[n] = \frac{X[n+1] - 2X[n] + X[n-1]}{\Delta t^2}, \quad a_y[n] = \frac{Y[n+1] - 2Y[n] + Y[n-1]}{\Delta t^2}, \quad a_z[n] = \frac{Z[n+1] - 2Z[n] + Z[n-1]}{\Delta t^2}.
\label{eq:central-acc}
\end{equation}
For boundary frames ($n=1$ and $n=N$), forward and backward finite difference approximations are utilized:
\begin{equation}
\vec{v}[1] = \frac{\vec{r}[2] - \vec{r}[1]}{\Delta t}, \quad \vec{v}[N] = \frac{\vec{r}[N] - \vec{r}[N-1]}{\Delta t}.
\end{equation}
The absolute linear velocity $v_{\text{abs}}[n]$ and absolute linear acceleration $a_{\text{abs}}[n]$ are calculated as Euclidean norms:
\begin{equation}
v_{\text{abs}}[n] = \sqrt{v_x[n]^2 + v_y[n]^2 + v_z[n]^2}, \quad a_{\text{abs}}[n] = \sqrt{a_x[n]^2 + a_y[n]^2 + a_z[n]^2}.
\end{equation}

\subsubsection{Joint Angle Determination via 3D Vector Geometry}
Joint kinematics are modeled through interconnected rigid segment vectors formed by marker pairs:
\begin{itemize}
    \item \textbf{Thigh Segment Vector ($\vec{u}_{\text{thigh}}$):} Directed from knee to hip, $\vec{u}_{\text{thigh}}[n] = \vec{r}_H[n] - \vec{r}_K[n]$.
    \item \textbf{Shank Segment Vector ($\vec{u}_{\text{shank}}$):} Directed from knee to ankle, $\vec{u}_{\text{shank}}[n] = \vec{r}_A[n] - \vec{r}_K[n]$.
    \item \textbf{Inverted Shank Vector ($\vec{w}_{\text{shank}}$):} Directed from ankle to knee, $\vec{w}_{\text{shank}}[n] = \vec{r}_K[n] - \vec{r}_A[n] = -\vec{u}_{\text{shank}}[n]$.
    \item \textbf{Foot Segment Vector ($\vec{u}_{\text{foot}}$):} Directed from ankle to foot marker, $\vec{u}_{\text{foot}}[n] = \vec{r}_F[n] - \vec{r}_A[n]$.
\end{itemize}

The three-dimensional knee flexion angle $\theta_{\text{knee}}[n]$ is derived using the vector dot product formula:
\begin{equation}
\theta_{\text{knee}}[n] = \arccos\left( \frac{\vec{u}_{\text{thigh}}[n] \cdot \vec{u}_{\text{shank}}[n]}{\|\vec{u}_{\text{thigh}}[n]\|\,\|\vec{u}_{\text{shank}}[n]\|} \right) \times \frac{180^\circ}{\pi}.
\label{eq:knee-angle}
\end{equation}
Under this standard three-point geometric formulation, full anatomical knee extension corresponds to $180^\circ$, whereas progressive knee flexion causes the angle to decrease below $180^\circ$ (flexion angle $\phi_{\text{knee}} = 180^\circ - \theta_{\text{knee}}$).

The ankle joint angle $\theta_{\text{ankle}}[n]$ (capturing dorsiflexion and plantarflexion) is computed between the inverted shank segment vector and the foot segment vector:
\begin{equation}
\theta_{\text{ankle}}[n] = \arccos\left( \frac{\vec{w}_{\text{shank}}[n] \cdot \vec{u}_{\text{foot}}[n]}{\|\vec{w}_{\text{shank}}[n]\|\,\|\vec{u}_{\text{foot}}[n]\|} \right) \times \frac{180^\circ}{\pi}.
\label{eq:ankle-angle}
\end{equation}
An angle of $90^\circ$ corresponds to neutral standing posture; angles exceeding $90^\circ$ indicate dorsiflexion, while angles below $90^\circ$ represent plantarflexion.

\subsection{Web-Based Simulation Environment (\csvfile{simulator.html})}
To evaluate 3D motion kinematics interactively without requiring commercial desktop software suites, an interactive web-based simulator was constructed using Three.js (WebGL) and Chart.js, contained in \csvfile{simulator.html}. The architecture features:
\begin{itemize}
    \item \textbf{Client-Side C3D Parser:} A custom JavaScript \texttt{C3DParser} class that ingests binary C3D files via \texttt{DataView} and \texttt{ArrayBuffer}, dynamically decoding header pointers, extracting parameter labels, and parsing floating-point frame records.
    \item \textbf{3D Scene Graph \& Skeletal Kinematic Rig:} Renders marker nodes as sphere meshes within a 3D coordinate frame featuring ground grids, directional lighting, and an orbit control camera. Skeletal segments connect anatomical joints dynamically based on proximal-distal parent relationships.
    \item \textbf{Dynamic Trajectory Trails:} Real-time historical trail line geometries illustrating the spatial trajectory curve traced by any selected marker.
    \item \textbf{Synchronized Telemetry Dashboard:} Interactive Chart.js canvas elements continuously plotting linear velocities and joint angles in synchrony with animation playback.
    \item \textbf{Interactive Playback Controls:} Full timeline scrubber, play/stop triggers, camera reset button, and adjustable playback speeds ($0.1\times$ to $3.0\times$).
\end{itemize}

\subsection{Dataset Summary}
The primary and secondary motion analysis datasets evaluated in this practical are detailed in Table~\ref{tab:datasets}.

\begin{table}[H]\centering
\caption{Experimental motion capture datasets and supporting simulation files.}
\begin{tabularx}{\textwidth}{llrrX}
\toprule
\textbf{Filename} & \textbf{Data Format} & \textbf{Size} & \textbf{Frames} & \textbf{Description} \\
\midrule
\csvfile{Group2_2026.c3d} & Binary C3D & $104.5\,\mathrm{kB}$ & 1644 & Primary lower-limb 3D motion capture trial ($100\,\mathrm{Hz}$) \\
\csvfile{Motion2026.c3d} & Binary C3D & $116.0\,\mathrm{kB}$ & 1820 & Supplementary motion capture trial for validation \\
\csvfile{Group2_2026_(Coordinates).txt} & ASCII Tabular & $907.4\,\mathrm{kB}$ & 1644 & 3D coordinates ($X, Y, Z$), linear velocities, accelerations, and residuals for 4 markers \\
\csvfile{Group2_2026.txt} & ASCII Tabular & $108.8\,\mathrm{kB}$ & 1644 & Computed knee and ankle joint angles, angular velocities, and angular accelerations \\
\csvfile{simulator.html} & HTML5 / WebGL & $41.4\,\mathrm{kB}$ & -- & Standalone 3D motion viewer, binary C3D parser, and kinematic telemetry dashboard \\
\bottomrule
\end{tabularx}
\label{tab:datasets}
\end{table}

\section{Results}

\subsection{Marker Trajectory Reconstruction}
The three-dimensional spatial coordinates ($X, Y, Z$) extracted from \csvfile{Group2_2026.c3d} and documented in \csvfile{Group2_2026_(Coordinates).txt} demonstrate distinct cyclical patterns corresponding to human walking gait cycles. Figure~\ref{fig:motion-trajectories} illustrates the reconstructed 3D spatial trajectories for the left hip, left knee, left ankle-bone, and left foot markers.

\begin{figure}[H]\centering
 \includegraphics[width=0.88\textwidth]{motion_3d_trajectories.png}
 \caption{Reconstructed three-dimensional spatial trajectories ($X, Y, Z$) of the left hip, left knee, left ankle-bone, and left foot markers across the capture trial.}
 \label{fig:motion-trajectories}
\end{figure}

The vertical ($Z$) displacement profiles exhibit clear anatomical stratification:
\begin{itemize}
    \item \textbf{Left Hip:} Displays minimal vertical translation, oscillating within a narrow band of $Z \approx 0.927\,\mathrm{m}$ to $0.943\,\mathrm{m}$ ($\Delta Z = 16\,\mathrm{mm}$). The hip marker demonstrates steady continuous forward displacement along the $+X$ progression axis, acting as the moving proximal reference frame.
    \item \textbf{Left Knee:} Exhibits coordinated vertical undulations between $Z \approx 0.485\,\mathrm{m}$ and $0.540\,\mathrm{m}$, lowering during stance phase flexion and elevating during swing phase initiation.
    \item \textbf{Left Ankle-Bone:} Alternates between stable stance support ($Z \approx 0.081\,\mathrm{m}$ to $0.095\,\mathrm{m}$) and swing phase peak elevation ($Z \approx 0.165\,\mathrm{m}$).
    \item \textbf{Left Foot:} Demonstrates flat ground contact during stance phase ($Z \approx 0.075\,\mathrm{m}$) followed by rapid elevation up to $Z \approx 0.203\,\mathrm{m}$ during mid-swing clearance.
\end{itemize}

A quantitative summary of the spatial displacement boundaries for each anatomical landmark is compiled in Table~\ref{tab:marker-displacements}.

\begin{table}[H]\centering
\caption{Spatial coordinate extremes and net displacement bounds for lower-limb markers.}
\begin{tabularx}{\textwidth}{lrrrrrr}
\toprule
\textbf{Marker} & \textbf{$X_{\min}$ (m)} & \textbf{$X_{\max}$ (m)} & \textbf{$Y_{\min}$ (m)} & \textbf{$Y_{\max}$ (m)} & \textbf{$Z_{\min}$ (m)} & \textbf{$Z_{\max}$ (m)} \\
\midrule
Left Hip & $0.005$ & $0.065$ & $-0.011$ & $0.026$ & $0.927$ & $0.943$ \\
Left Knee & $-0.036$ & $0.018$ & $0.097$ & $0.185$ & $0.485$ & $0.535$ \\
Left Ankle-Bone & $-0.087$ & $-0.069$ & $0.021$ & $0.065$ & $0.081$ & $0.106$ \\
Left Foot & $-0.097$ & $-0.091$ & $0.158$ & $0.203$ & $0.075$ & $0.095$ \\
\bottomrule
\end{tabularx}
\label{tab:marker-displacements}
\end{table}

\subsection{Kinematic Derivatives: Velocity and Acceleration Profiles}
Applying Equations~\eqref{eq:central-vel} and \eqref{eq:central-acc} yielded linear velocity and acceleration profiles across all three Cartesian axes. Figure~\ref{fig:marker-velocities} presents the absolute linear velocity curves for the four lower-limb markers.

\begin{figure}[H]\centering
 \includegraphics[width=0.88\textwidth]{marker_velocities.png}
 \caption{Absolute linear velocity profiles ($v_{\text{abs}}$) of the anatomical markers across the gait cycle, highlighting swing phase acceleration and stance phase deceleration.}
 \label{fig:marker-velocities}
\end{figure}

The velocity profiles capture key biomechanical events:
\begin{enumerate}
    \item \textbf{Stance Phase Quiescence:} During foot-flat stance (e.g., $t = 0.22\,\mathrm{s}$ to $0.35\,\mathrm{s}$), the left foot and ankle markers exhibit near-zero velocities ($v_{\text{abs}} < 0.02\,\mathrm{m/s}$), confirming rigid contact with the floor.
    \item \textbf{Swing Phase Surge:} During forward swing, the foot marker experiences rapid acceleration, reaching a peak linear velocity of $v_{\max} = 0.717\,\mathrm{m/s}$ at $t = 0.130\,\mathrm{s}$. The ankle-bone marker reaches a coincident peak velocity of $0.763\,\mathrm{m/s}$ at $t = 0.120\,\mathrm{s}$.
    \item \textbf{Hip Velocity Stability:} The hip marker moves at a relatively steady velocity fluctuating between $0.191\,\mathrm{m/s}$ and $0.751\,\mathrm{m/s}$, reflecting continuous whole-body forward momentum.
    \item \textbf{Impact Deceleration Spikes:} The acceleration traces reveal sharp deceleration spikes ($a_{\text{abs}} > 14.2\,\mathrm{m/s^2}$) at terminal swing immediately preceding initial contact, representing heel-strike impact transients.
\end{enumerate}

Table~\ref{tab:velocity-acceleration} lists the peak linear velocity, mean linear velocity, and peak absolute linear acceleration recorded for each anatomical marker.

\begin{table}[H]\centering
\caption{Kinematic derivative parameters for tracked anatomical landmarks.}
\begin{tabularx}{\textwidth}{lrrr}
\toprule
\textbf{Anatomical Marker} & \textbf{Peak Velocity $v_{\max}$ (m/s)} & \textbf{Mean Velocity $\bar{v}$ (m/s)} & \textbf{Peak Acceleration $a_{\max}$ ($\mathrm{m/s^2}$)} \\
\midrule
Left Hip & $0.751$ & $0.342$ & $6.552$ \\
Left Knee & $0.785$ & $0.298$ & $8.910$ \\
Left Ankle-Bone & $0.763$ & $0.264$ & $14.254$ \\
Left Foot & $0.717$ & $0.215$ & $12.599$ \\
\bottomrule
\end{tabularx}
\label{tab:velocity-acceleration}
\end{table}

\subsection{Joint Angle Analysis}
Joint angle trajectories computed from Equations~\eqref{eq:knee-angle} and \eqref{eq:ankle-angle} and documented in \csvfile{Group2_2026.txt} reveal characteristic sagittal-plane lower-extremity kinematics. Figure~\ref{fig:joint-angles} depicts the dynamic knee flexion-extension and ankle dorsiflexion-plantarflexion angle curves over time.

\begin{figure}[H]\centering
 \includegraphics[width=0.88\textwidth]{joint_angles.png}
 \caption{Dynamic joint angle trajectories for the left knee (flexion-extension) and left ankle (dorsiflexion-plantarflexion) across consecutive gait cycles.}
 \label{fig:joint-angles}
\end{figure}

Key observations from the joint kinematics include:
\begin{enumerate}
    \item \textbf{Knee Joint Kinematics:} The knee angle ranges from a minimum of $151.31^\circ$ (maximum flexion during stance/swing transition) to a maximum of $169.74^\circ$ (near full extension prior to heel strike), defining an angular range of motion of $\mathrm{ROM}_{\text{knee}} = 18.43^\circ$. Maximum knee extension velocity reached $\omega_{\text{ext}} = 401.84^\circ/\mathrm{s}$ at $t = 1.68\,\mathrm{s}$, driven by rapid quadriceps concentric contraction during terminal swing.
    \item \textbf{Ankle Joint Kinematics:} The ankle angle oscillates between a minimum of $76.08^\circ$ (plantarflexion during push-off at $t = 2.13\,\mathrm{s}$) and a maximum of $99.70^\circ$ (dorsiflexion during mid-swing clearance at $t = 1.82\,\mathrm{s}$), producing an ankle range of motion of $\mathrm{ROM}_{\text{ankle}} = 23.62^\circ$. Peak ankle angular acceleration reached $\alpha_{\max} = 19862.2^\circ/\mathrm{s^2}$ at $t = 1.78\,\mathrm{s}$, corresponding to rapid triceps surae elastic recoil and toe push-off.
\end{enumerate}

Table~\ref{tab:joint-kinematics} details the quantitative kinematic metrics for the knee and ankle joints.

\begin{table}[H]\centering
\caption{Lower-limb joint angular kinematics summary.}
\begin{tabularx}{\textwidth}{lrrrrr}
\toprule
\textbf{Joint Angle} & \textbf{Min ($^\circ$)} & \textbf{Max ($^\circ$)} & \textbf{ROM ($^\circ$)} & \textbf{Peak Velocity ($^\circ$/s)} & \textbf{Peak Acceleration ($^\circ/\mathrm{s^2}$)} \\
\midrule
Left Knee & $151.31$ & $169.74$ & $18.43$ & $401.84$ & $7274.19$ \\
Left Ankle & $76.08$ & $99.70$ & $23.62$ & $298.56$ & $19862.22$ \\
\bottomrule
\end{tabularx}
\label{tab:joint-kinematics}
\end{table}

\subsection{Interactive Web Simulator Verification}
Loading \csvfile{Group2_2026.c3d} into \csvfile{simulator.html} validated the computational pipeline in real time. The custom Three.js engine confirmed:
\begin{enumerate}
    \item \textbf{Segment Length Invariance:} Distances between anatomical marker pairs remained structurally invariant throughout animation ($L_{\text{thigh}} \approx 0.44\,\mathrm{m}$, $L_{\text{shank}} \approx 0.42\,\mathrm{m}$, $L_{\text{foot}} \approx 0.15\,\mathrm{m}$), demonstrating rigid body fidelity with negligible soft tissue distortion.
    \item \textbf{Telemetry Synchronization:} Real-time dynamic updates in Chart.js verified instantaneous agreement between 3D marker velocity vectors and displayed telemetry values.
    \item \textbf{Render Performance:} Stable execution at 60 frames per second (FPS) with smooth timeline scrubbing across all 1644 frames.
\end{enumerate}

\section{Discussion}

\subsection{Integrity and Robustness of C3D Binary Parsing}
Manual and programmatic extraction of binary C3D structures confirms the exceptional utility of the format for biomechanical research. By storing data in standardized 512-byte blocks with strict byte-offset pointers, C3D eliminates the parsing ambiguities, localization errors, and delimiter inconsistencies prevalent in text-based tabular formats (e.g., CSV or TSV). In particular, reading the header processor type (\texttt{0x50}) guarantees correct little-endian IEEE 754 decoding on modern computing platforms. Furthermore, embedding the floating-point scale factor (\texttt{POINT:SCALE}) and camera residual thresholds directly within parameter records enables automated validation of coordinate tracking accuracy without external calibration logs.

\subsection{Noise Amplification in Kinematic Differentiation}
A recognized challenge in optical motion capture analysis is high-frequency noise amplification during numerical differentiation. As demonstrated by Equation~\eqref{eq:central-acc}, differentiating discrete displacement data twice via central finite differences introduces a second-order noise sensitivity proportional to $1/\Delta t^2 = 1/(0.01)^2 = 10^4$. Consequently, minor optical measurement jitter of $\pm 0.5\,\mathrm{mm}$ induces substantial acceleration fluctuations exceeding $5000^\circ/\mathrm{s^2}$ (as observed in the raw ankle acceleration data where peaks reach $\sim 19862^\circ/\mathrm{s^2}$). 

In clinical gait analysis routines, this artefact is mitigated by applying a zero-lag fourth-order low-pass Butterworth filter or quintic smoothing splines with an optimal cut-off frequency between $6\,\mathrm{Hz}$ and $10\,\mathrm{Hz}$. Filtering attenuates high-frequency camera digitization jitter while preserving genuine physiological accelerations.

\subsection{Soft Tissue Artefact and Marker Dropouts}
Optical tracking assumes that surface-mounted retroreflective markers remain rigidly attached to the underlying skeletal framework. In practice, soft tissue artefact (STA)---caused by muscle deformation, skin stretching, and inertial wobble relative to bone---introduces localized kinematic errors ranging from $5\,\mathrm{mm}$ to $20\,\mathrm{mm}$, particularly at the thigh and pelvis. In our dataset, inspection of early frames ($t = 0.00\,\mathrm{s}$ to $0.06\,\mathrm{s}$) reveals temporary marker unassigned periods (#2D camera residual flags), requiring piecewise cubic Hermite or linear interpolation to reconstruct continuous kinematic series.

\subsection{Utility of Web-Based Motion Simulation}
The implementation of \csvfile{simulator.html} highlights the growing viability of browser-based WebGL environments for clinical and educational biomechanics. Traditional motion capture analysis relies on proprietary commercial packages (such as Vicon Nexus, Qualisys Track Manager, or Motion Analysis Cortex) that demand high-end dedicated workstations and expensive licensing. 

In contrast, our Three.js simulator runs entirely client-side on any modern web browser without server-side dependencies or software installation. By combining real-time 3D skeletal rendering with dynamic Chart.js telemetry, clinicians and students can visualize gait kinematics, inspect marker velocity vectors, and inspect joint angles synchronously. This democratization of motion capture data facilitates remote clinical consultations, student laboratory learning, and telemedicine evaluations.

\section{Conclusion}

This practical successfully investigated the complete three-dimensional optical motion capture processing pipeline. The binary C3D file format was parsed to extract anatomical marker coordinates, frame metadata, and acquisition parameters at $100\,\mathrm{Hz}$. Three-dimensional trajectories of the left hip, left knee, left ankle-bone, and left foot were reconstructed across full gait cycles. Linear velocities and accelerations were derived using central finite difference approximations, capturing swing phase accelerations and heel-strike deceleration transients. Knee flexion-extension and ankle dorsiflexion-plantarflexion angles were computed using 3D vector dot product formulations, revealing characteristic joint ranges of motion of $18.43^\circ$ and $23.62^\circ$, respectively. Finally, an interactive Three.js web simulator (\csvfile{simulator.html}) was deployed to provide real-time 3D skeletal rendering and synchronized kinematic telemetry.

\end{document}
